% !TEX root = ../main.tex
% Section 1. Page budget: about 0.5 page (trimming pass, 2026-11-24).
% Sources: PROJECT_CONTEXT.md §1 (summary), §2 (deviations), §5 (improvements); course brief §5
% (Option E wording). D-numbers: D6 (pure P2P), D10 (CDC), D11 (tagged discovery).
\vspace{0.6pt}
\section{Introduction and the option selected}
\label{sec:intro}

meshsync keeps one shared folder identical on three Linux hosts, with no central server: every
host is an equal peer, and a change on one reaches the others within seconds. We chose a custom
option extending Option~E (file distribution): a peer-to-peer synchroniser with our own protocol on
sockets, UDP broadcast for discovery and TCP for everything else. Instead of a scaling study,
Section~\ref{sec:ftp} compares it with FTP, a client-server protocol, on the same hosts. Each
feature fixes a gap in an existing protocol: only changed chunks travel, unlike FTP, even after an
insert (content-defined chunking); request IDs allow out-of-order replies, unlike HTTP/1.1; many peers
stay in sync continuously, unlike rsync; and an HMAC handshake never sends the secret, unlike FTP's
clear-text password.

\paragraph{Approved deviations.} From our approved proposal, quoted verbatim:

\begin{quote}
Approved deviations (stated in proposal): custom option (extends Option E); R5 capture in plain
mode; no client-server vs P2P scaling comparison.
\end{quote}
